\section{Conclusion and Future Work}
\label{sec:conclusion}

We have presented, to our knowledge, the first kernel-checked proof, with no
project-specific axioms, of norm convergence of the Lie--Trotter sequence for
arbitrary elements of a complete normed algebra, together with higher-order
extensions and commutator-scaling error bounds, in Lean~4 with the \Mathlib
library.
The development comprises over 14\,000 lines with zero \texttt{sorry} axioms
and zero own theorem-level BCH-interface axioms. Levels~1--4 are fully
proved at the project level: at Lean-BCH revision \texttt{05e8c52},
\texttt{\#print axioms} returns only \texttt{[propext,
Classical.choice, Quot.sound]} on both the $\tau^5$ headlines and the
Level~4 local small-$\tau$ ($\tau^7$) refinement. The development covers:
\begin{itemize}
  \item First-order Lie--Trotter convergence at rate $O(1/n)$;
  \item Symmetric Strang splitting at rate $O(1/n^2)$;
  \item Multi-operator generalizations of both formulas;
  \item First- and second-order commutator-scaling bounds via a
    Duhamel/FTC technique;
  \item A second-order Strang bound via norm-of-difference that captures
    signed cancellations invisible to the sum-of-norms form, and whose
    leading coefficient is machine-checked never to exceed it---strictly
    smaller whenever the two double commutators partially cancel, a gain
    measured at $9$--$50\%$ across standard spin-chain benchmarks
    (\cref{apd:tighter});
  \item A fourth-order Suzuki $S_4$ track with four CAS-assisted
    BCH-based bounds, including a replacement of the 2021 prefactors of
    Childs et al.\ (rigorous, but not proven tight) by CAS-certified
    coefficients smaller term by term by 8.6$\times$--64$\times$---the
    comparison machine-checked at the leading coefficient, with full-bound
    domination on a neighbourhood of zero under a strict-gap
    hypothesis---and an existential-$\delta$ small-$\tau$ refinement via
    an explicit $\tau^7$ correction;
  \item Fourth-order Suzuki convergence at rate $O(1/n^4)$, obtained by
    compounding the single-step $O(t^5)$ bound over $n$ steps --- and, from
    the sharp step bound, with an explicit commutator-scaled leading
    coefficient.
\end{itemize}
The first, second, and fourth-order convergence theorems together form a
complete verified hierarchy $O(1/n) \to O(1/n^2) \to O(1/n^4)$: for each
integrator the development supplies both a single-step error bound and the
compounded statement that the $n$-step product converges to $e^{t(A+B)}$ at the
corresponding rate.

The Duhamel approach to commutator scaling---avoiding ODE uniqueness in favor
of the fundamental theorem of calculus---may be of independent interest as a
proof strategy for related product formula bounds.

\paragraph{Future directions.}
Several extensions are natural:
\begin{enumerate}
  \item \emph{Mathlib contributions.}
    The exponential norm bounds (\leanref{ExpBounds.lean}{72}{norm\_exp\_le},
    \leanref{ExpBounds.lean}{88}{norm\_exp\_sub\_one\_le},
    \leanref{ExpBounds.lean}{129}{exp\_sub\_one\_sub\_bound\_real})
    are general-purpose lemmas suitable for upstreaming to \Mathlib.

  \item \emph{Trotter-native and higher-order BCH derivations.}
    The three-factor symmetric BCH identities, the two-factor
    palindromic quintic remainder, and the five-factor palindromic
    quintic and septic identifications are formalized in the Lean-BCH
    companion project at revision~\texttt{05e8c52}. A Trotter-native
    operator-algebra proof of the order-4 Taylor-coefficient identity
    (via \texttt{scripts/compute\_sumQuadCorr\_factored.py}) remains a
    useful independent derivation, while extending the certified BCH
    expansion to order nine would sharpen the local remainder.

  \item \emph{Specialization to quantum systems.}
    Connecting our abstract bounds to concrete gate-count estimates for
    Hamiltonian simulation on quantum computers would bridge the gap between
    formal mathematics and verified quantum algorithms.

  \item \emph{Unbounded operators.}
    The original Trotter--Kato theorem applies to strongly continuous
    semigroups on Banach spaces, covering unbounded generators.
    Formalizing this extension would require substantial additional
    functional analysis infrastructure in \Mathlib.
\end{enumerate}

More broadly, our project demonstrates that substantial results in operator
analysis---involving non-commutative algebra, power series manipulation,
Fr\'{e}chet calculus, and Bochner integration---are within reach of current
proof assistant technology.
We hope this work encourages further formalization of the mathematical
foundations underlying quantum algorithms.
